% !TeX root = ../memoria.tex
\chapter{Introducción}

    Se disposne del código fuente, ejemplos y procedimientos de instalación en un 
    \href{https://github.com/andrelorenzo/TFM_2dvision}{repositiorio} público de GitHub.
\section{Motivación}
    Uno de los principales retos de la navegación autónoma en UAV es la selección y el uso de sensores para la evitación de obstáculos. 
    Debido a las restricciones de peso y volumen, en la mayoría de los casos resulta inviable emplear LiDAR 3D o cámaras RGB-D. Además, 
    existe una limitación adicional de carácter fundamental: este tipo de sensores genera datos tridimensionales cuyo procesamiento 
    requiere una capacidad de cómputo elevada, así como unas prestaciones que, por lo general, no están disponibles en drones de pequeño 
    tamaño.

    Durante la realización del trabajo tutelado de investigación se llevó a cabo un estudio en profundidad sobre los métodos empleados 
    para lograr una navegación segura y una evitación de obstáculos eficaz~\cite{estudio-nav2D}. Asimismo, se analizaron múltiples modelos 
    de estimación de profundidad, comparándolos en términos de requisitos computacionales, velocidad de inferencia, estabilidad en la 
    detección y precisión de la medida. Los buenos resultados obtenidos en el citado trabajo de investigación han sido una gran fuente 
    de motivación para la realización de este trabajo.

    A diferencia de otros vehículos autónomos, en los que se dispone de sensórica y de unidades de procesamiento capaces de estimar la 
    posición con alta precisión, en este entorno la limitada disponibilidad de datos obliga a desarrollar soluciones optimizadas para la 
    navegación autónoma. Este condicionante requiere priorizar arquitecturas ligeras y robustas, capaces de operar en tiempo real con 
    recursos computacionales restringidos, sin comprometer los márgenes de seguridad durante el vuelo.

    En consecuencia, el planteamiento adoptado en este trabajo se orienta hacia el uso de sensores y métodos de percepción que mantengan 
    un equilibrio adecuado entre coste computacional, fiabilidad y precisión. En particular, se propone explotar información visual 
    procedente de cámaras monoculares, complementada con modelos de estimación de profundidad eficientes, con el objetivo de generar 
    una representación útil del entorno para la detección de obstáculos y la planificación de trayectorias.

\section{Objetivos del estudio}
    El objetivo principal de este trabajo es dotar a un UAV de pequeño tamaño, con prestaciones limitadas, de la capacidad de navegar de 
    forma autónoma en entornos interiores. Para ello, el vehículo deberá ser capaz de desplazarse siguiendo una trayectoria previamente 
    definida por el usuario, a la vez que evita los obstáculos presentes en su recorrido. A continuación, se enumeran los requisitos 
    fundamentales cuya consecución permitirá considerar satisfactoria la finalización del trabajo.

    \begin{itemize}
        \item \textbf{Sistema de sensorización:} El conjunto de sensores queda estipulado y limitado, de antemano, a una cámara monocular 
        y una unidad de medida inercial (IMU).
        \item \textbf{Tiempo real:} El conjunto de algoritmos deberá ejecutarse dentro de un límite temporal de 100~ms, equivalente a una 
        frecuencia de 10~Hz, con el fin de considerar el sistema utilizable en aplicaciones en tiempo real.
        \item \textbf{Ordenador central:} El ordenador central realizará la mayor parte del cómputo, excluyendo la captura de imágenes y un 
        preprocesamiento básico destinado a mejorar la calidad o reducir el tamaño de las imágenes obtenidas.
        \item \textbf{Ground truth:} Los resultados se validarán mediante información tridimensional capturada por la cámara de abordo. Sin 
        embargo, dicha información no podrá emplearse en el algoritmo de navegación, utilizándose exclusivamente como \textit{ground truth}. 
        Esta información se almacenará en la memoria a bordo y no se transmitirá al sistema de cómputo central.
        \item \textbf{Precisión:} La precisión del sistema se medirá como la diferencia entre la trayectoria estipulada por el usuario y la 
        trayectoria real del vehículo. En ningún tramo del recorrido dicha diferencia deberá superar los 0.5~m.
        \item \textbf{Seguridad:} El UAV deberá ser capaz de esquivar obstáculos estáticos y dinámicos manteniendo una distancia de 
        seguridad equivalente al doble de su \textit{bounding box}.
        \item \textbf{Rango de operabilidad:} Dado que el objetivo principal del trabajo es el estudio de la navegación y no de las 
        comunicaciones, el alcance máximo de operación del algoritmo estará limitado por la tecnología inalámbrica empleada. Este alcance 
        se definirá como un radio centrado en la posición del ordenador central.
        \item \textbf{Comunicación:} La comunicación con el vehículo no constituye un objeto de estudio en sí mismo, al considerarse 
        independiente del algoritmo propuesto. En este trabajo se emplearán comunicaciones UDP y RTSP para el envío y la recepción de datos.
        \item \textbf{Optimizaciones:} Dada la relevancia de la capacidad de cómputo en el sistema, las optimizaciones se circunscribirán 
        al hardware específico empleado. No obstante, se aportará información suficiente para permitir la replicación de los resultados en 
        otros dispositivos. El ordenador de a bordo será una Jetson Nano (reComputer J1010) y el ordenador central un portátil HP ProBook 
        con procesador Intel Core Ultra 7 y GPU NVIDIA GeForce RTX 2050. La cámara será una Realsense 530.
    \end{itemize}
\clearpage
\section{Estructura del trabajo}
    Este trabajo se estructura en nueve capítulos:

    En el primer capítulo, correspondiente a la presente introducción, se exponen las motivaciones, los requisitos y los objetivos del 
    proyecto, así como la definición de la estructura general del documento.
    
    El segundo capítulo presenta el estado del arte en navegación basada en visión monocular y en navegación monocular combinada con 
    información inercial, incluyendo las técnicas más empleadas y representativas en el momento de realización del proyecto.
    
    En el tercer capítulo se define la arquitectura técnica del sistema, incorporando las principales decisiones de diseño, el esquema de 
    comunicaciones y las distintas etapas de procesado, tanto a nivel de sistema como en el procesamiento de imágenes.
    
    El cuarto capítulo se centra en el estudio de viabilidad del uso de modelos de estimación de profundidad~\cite{estudio-nav2D}, 
    describiendo los modelos seleccionados, su arquitectura, las optimizaciones aplicadas y las métricas empleadas para su evaluación.
    
    El quinto capítulo aborda la implementación del controlador de navegación basado en el estimador de profundidad.
    
    El sexto capítulo describe el desarrollo e integración de una solución VI-SLAM optimizada para el contexto del proyecto, detallando 
    su utilización como planificador global y/o local.
    
    El séptimo capítulo recoge las principales dificultades encontradas durante el desarrollo, así como las decisiones de diseño revisadas 
    para alcanzar una solución viable.
    
    El octavo capítulo presenta los ensayos y pruebas realizados para validar la solución propuesta.
    
    Finalmente, el noveno capítulo expone las conclusiones obtenidas y las líneas de trabajo futuro.

